\documentclass[11pt]{article}
\usepackage{amsmath,amssymb,amsthm}
\usepackage[margin=1in]{geometry}

\newtheorem{theorem}{Theorem}
\newtheorem{lemma}[theorem]{Lemma}

\title{Problem 525: Rolling Ellipse}
\author{Project Euler Solutions}
\date{}

\begin{document}
\maketitle

\section{Problem Statement}
An ellipse with semi-major axis~$a$ and semi-minor axis~$b$ ($a > b > 0$) rolls without slipping along the $x$-axis. Compute the arc length of the path traced by the center of the ellipse over one complete revolution.

\section{Mathematical Foundation}

\begin{theorem}[Ellipse Arc Length]
The arc length of the ellipse $x^2/a^2 + y^2/b^2 = 1$ from parameter $0$ to $\theta$ is
\[
s(\theta) = a\,E(\theta \mid e^2)
\]
where $e = \sqrt{1 - b^2/a^2}$ is the eccentricity and $E(\phi \mid m)$ is the incomplete elliptic integral of the second kind.
\end{theorem}

\begin{proof}
Parameterize the ellipse as $\mathbf{r}(t) = (a\cos t,\, b\sin t)$. Then
\[
|\mathbf{r}'(t)| = \sqrt{a^2\sin^2 t + b^2\cos^2 t} = a\sqrt{1 - e^2\cos^2 t}.
\]
Integrating from $0$ to $\theta$:
\[
s(\theta) = a\int_0^{\theta}\!\sqrt{1 - e^2\cos^2 t}\,dt = a\,E(\theta \mid e^2). \qedhere
\]
\end{proof}

\begin{lemma}[Distance from Center to Tangent]
The perpendicular distance from the center of the ellipse to the tangent line at parameter~$\theta$ is
\[
h(\theta) = \frac{ab}{\sqrt{a^2\sin^2\theta + b^2\cos^2\theta}}.
\]
\end{lemma}

\begin{proof}
The tangent line at $(a\cos\theta, b\sin\theta)$ is $\frac{x\cos\theta}{a} + \frac{y\sin\theta}{b} = 1$. The distance from the origin to this line is
\[
h = \frac{1}{\sqrt{\cos^2\theta/a^2 + \sin^2\theta/b^2}} = \frac{ab}{\sqrt{b^2\cos^2\theta + a^2\sin^2\theta}}. \qedhere
\]
\end{proof}

\begin{theorem}[Center Trajectory]
When the ellipse has rotated so that the contact point is at parameter~$\theta$, the center is at
\[
X_c(\theta) = s(\theta) - \frac{(a^2-b^2)\sin\theta\cos\theta}{\sqrt{a^2\sin^2\theta + b^2\cos^2\theta}}, \qquad
Y_c(\theta) = \frac{ab}{\sqrt{a^2\sin^2\theta + b^2\cos^2\theta}}.
\]
\end{theorem}

\begin{proof}
The rolling constraint equates arc length along the ellipse with distance along the $x$-axis: $x_{\mathrm{contact}} = s(\theta)$. The center is displaced from the contact point by the perpendicular distance~$h(\theta)$ along the inward normal. Decomposing into horizontal and vertical components using the tangent angle yields the stated formulas.
\end{proof}

\begin{theorem}[Arc Length of Center Path]
The arc length of the center's trajectory over one complete revolution is
\[
L = \int_0^{2\pi}\!\sqrt{\Bigl(\frac{dX_c}{d\theta}\Bigr)^{\!2} + \Bigl(\frac{dY_c}{d\theta}\Bigr)^{\!2}}\,d\theta.
\]
This integral does not admit a closed form in terms of standard functions.
\end{theorem}

\begin{proof}
The arc-length formula for a parametric curve is standard. The integrand involves compositions of trigonometric functions and square roots of rational expressions in $\sin\theta,\cos\theta$ with parameters $a,b$, which cannot be reduced to elementary or standard elliptic integrals.
\end{proof}

\section{Algorithm}

\begin{enumerate}
    \item Compute the derivatives $dX_c/d\theta$ and $dY_c/d\theta$ symbolically.
    \item Evaluate the arc-length integral $L$ numerically using adaptive Gauss--Kronrod quadrature with the required precision.
\end{enumerate}

\section{Complexity Analysis}
\begin{itemize}
    \item \textbf{Time:} $O(N)$ for $N$ quadrature points; adaptive methods achieve arbitrary precision.
    \item \textbf{Space:} $O(N)$ for the quadrature recursion stack; $O(1)$ for fixed-rule quadrature.
\end{itemize}

\section{Answer}

\[
\boxed{44.69921807}
\]

\end{document}
